% Bagian: first-order-logic
% Bab: introduction
% Subbagian: formulas

\documentclass[../../../include/open-logic-section]{subfiles}

\begin{document}

\olfileid[id]{fol}{int}{fml}

\section{\usetoken{P}{formula}}

Berikut adalah pendekatan yang akan kita gunakan untuk menetapkan secara
ketat !!{sentence}s logika orde pertama dan menangani persoalan yang timbul
dari penggunaan !!{variable}s. Mula-mula kita mendefinisikan suatu himpunan
ungkapan yang \emph{berbeda}: !!{formula}s. Setelah itu dilakukan, kita dapat
mempertimbangkan peran !!{variable}s di dalamnya---dan berdasarkan beberapa
gagasan lain, yakni !!{variable}s ``bebas'' dan ``terikat,'' kita dapat
mendefinisikan apa itu !!a{sentence} (yakni, !!a{formula} tanpa
!!{variable}s bebas). Kita melakukan ini bukan hanya karena cara tersebut
membuat definisi ``!!{sentence}'' lebih mudah ditangani, melainkan juga karena
hal ini sangat penting bagi cara kita mendefinisikan gagasan semantis
keterpenuhan.

Mari kita definisikan ``!!{formula}'' bagi suatu bahasa orde pertama
sederhana, yang hanya memuat satu !!{predicate}~$\Obj P$ dan satu
!!{constant}~$\Obj a$, serta hanya simbol logis $\lnot$, $\land$,
dan~$\lexists$. Definisi lengkap kita akan jauh lebih umum: kita akan
mengizinkan tak hingga banyaknya !!{predicate}s dan !!{constant}s. Bahkan,
kita juga akan mempertimbangkan !!{function}s yang dapat digabungkan dengan
!!{constant}s dan !!{variable}s untuk membentuk ``suku.'' Untuk sekarang,
$\Obj a$ dan variabel-variabel akan menjadi satu-satunya suku kita. Kita
memang memerlukan tak hingga banyaknya !!{variable}s. Secara resmi, kita akan
menggunakan simbol $\Obj v_0$, $\Obj v_1$, \dots, sebagai variabel.

\begin{defn}
Himpunan \emph{!!{formula}s}~$\Frm$ didefinisikan sebagai berikut:
\begin{enumerate}
\item\ollabel{fmls-atom} $\Atom{\Obj P}{\Obj a}$ dan $\Atom{\Obj
  P}{\Obj v_i}$ merupakan !!{formula}s ($i \in \Nat$).

\tagitem{prvNot}{\ollabel{fmls-not}Jika $!A$ merupakan !!a{formula}, maka
  $\lnot !A$ merupakan !!{formula}.}{}

\tagitem{prvAnd}{Jika $!A$ dan $!B$ merupakan !!{formula}s, maka $(!A \land
  !B)$ merupakan !!a{formula}.}{}

\tagitem{prvEx}{\ollabel{fmls-ex}Jika $!A$ merupakan !!a{formula} dan $x$
  merupakan !!a{variable}, maka $\lexists[x][!A]$ merupakan
  !!a{formula}.}{}

\tagitem{limitClause}{\ollabel{fmls-limit}Tidak ada ungkapan lain yang
  merupakan !!a{formula}.}{}
\end{enumerate}
\end{defn}

\olref{fmls-atom} memberi tahu kita bahwa $\Atom{\Obj P}{\Obj a}$ dan
$\Atom{\Obj P}{\Obj v_i}$ merupakan !!{formula}s untuk setiap $i \in
\Nat$. Semua ini disebut !!{formula}s \emph{atomik}. Formula-formula tersebut
memberi kita titik awal. Kasus-kasus lain memberi kita cara untuk membentuk
!!{formula}s baru dari formula yang telah kita bentuk. Jadi, misalnya,
berdasarkan \olref{fmls-not}, kita memperoleh bahwa $\lnot \Atom{\Obj P}{\Obj
v_2}$ merupakan !!a{formula}, karena $\Atom{\Obj P}{\Obj v_2}$ sudah
merupakan !!a{formula} berdasarkan \olref{fmls-atom}. Lalu, berdasarkan
\olref{fmls-ex}, kita memperoleh bahwa $\lexists[\Obj v_2][\lnot
\Atom{\Obj P}{\Obj v_2}]$ merupakan !!{formula} lain, dan demikian
seterusnya. \olref{fmls-limit} memberi tahu kita bahwa \emph{hanya} untai
yang dapat kita bentuk dengan cara inilah yang tergolong !!{formula}s. Secara
khusus, $\lexists[\Obj v_0][\Atom{\Obj P}{\Obj a}]$ dan $\lexists[\Obj
v_0][\lexists[\Obj v_0][\Atom{\Obj P}{\Obj a}]]$ \emph{memang} tergolong
!!{formula}s, sedangkan $(\lnot \Atom{\Obj P}{\Obj a})$ tidak, karena adanya
pasangan tanda kurung luar yang tidak perlu.

Cara mendefinisikan !!{formula}s ini disebut \emph{definisi induktif}, dan
cara ini memungkinkan kita membuktikan berbagai hal mengenai !!{formula}s
dengan menggunakan suatu versi bukti melalui induksi yang disebut
\emph{induksi struktural}. Hal-hal ini dibahas secara umum dalam
\olref[mth][ind][idf]{sec} dan \olref[mth][ind][sti]{sec}, yang sebaiknya Anda
tinjau sebelum mendalami bukti-bukti selanjutnya. Pada dasarnya, gagasannya
ialah bahwa jika Anda hendak membuktikan bahwa sesuatu benar bagi semua
!!{formula}s, pertama-tama Anda menunjukkan bahwa hal itu benar bagi
!!{formula}s atomik, lalu menunjukkan bahwa \emph{jika} hal itu benar bagi
sebarang !!{formula}~$!A$ (dan~$!B$), hal itu \emph{juga} benar bagi $\lnot
!A$, $(!A \land !B)$, dan $\lexists[x][!A]$. Misalnya, cara ini membuktikan
bahwa hal tersebut benar bagi $\lexists[\Obj v_2][\lnot \Atom{\Obj P}{\Obj
v_2}]$: dari bagian pertama Anda mengetahui bahwa hal tersebut benar bagi
!!{formula} atomik~$\Atom{\Obj P}{\Obj v_2}$. Kemudian, dari bagian kedua
Anda memperoleh bahwa hal tersebut benar bagi $\lnot \Atom{\Obj P}{\Obj
v_2}$, lalu sekali lagi bahwa hal tersebut benar bagi $\lexists[\Obj
v_2][\lnot \Atom{\Obj P}{\Obj v_2}]$ itu sendiri. Karena semua !!{formula}s
dibangkitkan secara induktif dari !!{formula}s atomik, cara ini berlaku bagi
formula mana pun.

\end{document}
